\begin{figure}[!t]\centering
\includegraphics[width=\linewidth,height=0.82\textheight,keepaspectratio]{figures/Figure_1.pdf}
\caption{\textbf{无沟通条件下的贡献动态。} 各面板展示六个推理配置在N=6、R=0.40、0.90、0.98和1.10时的结果。点表示各轮群体贡献率的系综均值，阴影表示在条件内对完整轨迹进行10,000次bootstrap重抽样得到的逐点95\%百分位区间。全部博弈采用T=10、每轮重新发放禀赋E=20、固定成员、汇总反馈和全部历史。}
\label{fig:1}
\end{figure}

\begin{figure}[!t]\centering
\includegraphics[width=\linewidth,height=0.82\textheight,keepaspectratio]{figures/Figure_2.pdf}
\caption{\textbf{固定群体规模下随激励变化的贡献水平。} 全部面板采用N=6和无沟通条件。（a）全部11个R值下第1--10轮的全过程平均贡献率。（b）相同扫描范围内第10轮的终轮贡献率。（c）将a中的六个社会困境R值单独展开。R=r/N按实际数值间距绘制，a和b中的竖虚线标记R=1。点和阴影分别为条件均值及逐点95\%完整轨迹bootstrap区间，重抽样10,000次。其他设置与图1相同。}
\label{fig:2}
\end{figure}

\begin{figure}[!t]\centering
\includegraphics[width=\linewidth,height=0.82\textheight,keepaspectratio]{figures/Figure_3.pdf}
\caption{\textbf{公开沟通与贡献动态。} 每个模型区域包含R=0.60、0.98、1.00和1.02下的四个时间序列面板（a--d），以及全过程沟通差值图（e）。灰色圆点、蓝色圆点和橙色方块分别表示无沟通、首轮前公开沟通和每轮前公开沟通。在e中，CE表示沟通效应（communication effect），pp是百分点（percentage points）的缩写。CE为沟通条件下的全过程贡献率均值减去相应无沟通均值，再乘以100。First和Every分别标识两种沟通安排。五行对应R=0.60、0.90、0.98、1.00和1.02。单元格数值四舍五入至整数百分点，精确差值及区间保留于源数据。均值及差值基于各自独立的完整系综计算。时间曲线的阴影为在条件内对完整轨迹进行10,000次bootstrap重抽样得到的逐点95\%区间。全部博弈采用N=6、T=10、E=20、固定成员、全部历史和汇总反馈。}
\label{fig:3}
\end{figure}

\begin{figure}[!t]\centering
\includegraphics[width=\linewidth,height=0.82\textheight,keepaspectratio]{figures/Figure_4.pdf}
\caption{\textbf{不同轮次与沟通安排下的贡献行动构成。} 每个模型区域依次展示无沟通、首轮前公开沟通和每轮前公开沟通。N=6时，先计算条件内比例，再对R=0.60、0.90和0.98等权平均。零贡献、部分贡献和全额贡献分别对应\(c_i(t)=0\)、\(0<c_i(t)<E\)和\(c_i(t)=E\)。每个实际观察轮次的各类比例之和为100\%。面积连接十个轮次的取值，不显示不确定性区间。每个条件包含的独立群体数在DeepSeek中为10--13，其余模型均为30。}
\label{fig:4}
\end{figure}

\begin{figure}[!t]\centering
\includegraphics[width=\linewidth,height=0.82\textheight,keepaspectratio]{figures/Figure_5.pdf}
\caption{\textbf{随机终止与固定轮数的比较。} 上部面板展示无沟通条件下R=0.90和1.02、\ensuremath{\delta}=0.90、N=6、E=20下各个已完成群体的轨迹。每行对应一个群体，按实际轮数排序。彩色单元格表示实际观察的逐轮贡献率，黑色短线标记实际终点。灰色斜线区域表示终止后未观察的时段。下部面板展示固定十轮或随机终止时各群体的全过程平均贡献率。水平抖动仅用于分开同一类别中的点。黑色菱形及误差棒表示群体等权均值，以及对完整轨迹进行10,000次bootstrap重抽样得到的95\%区间。停止规则见补充材料S1.5节。}
\label{fig:5}
\end{figure}

\begin{figure}[!t]
\centering
\includegraphics[width=\linewidth,height=0.78\textheight,keepaspectratio]{figures/Figure_6.pdf}
\caption{\textbf{群体规模与公开沟通。} 面板顺序与图1一致。灰色圆点、蓝色圆点和橙色方块分别表示无沟通、首轮公开沟通和每轮公开沟通。各点为N=4、6、8或12下的全过程均值，对R=0.60、0.90和0.98等权平均。误差棒为在条件内对完整群体轨迹进行10,000次bootstrap重抽样得到的逐点95\%百分位区间。比较包含216个条件和3,977个完整群体，每条件9—30组。N=4、8和12下的沟通条件各有10组，其中三个Kimi条件有九个已完成群体。全部博弈采用T=10、E=20、预先告知终点、固定成员、全部历史和汇总反馈。N改变时R保持固定。虚线连接各规模下的估计值。}
\label{fig:6}
\end{figure}

\begin{figure}[H]
\centering
\includegraphics[width=\linewidth,height=0.78\textheight,keepaspectratio]{figures/Figure_7.pdf}
\caption{\textbf{不同激励与群体规模下的推理设置比较。} DeepSeek对应a1--d1，Claude对应a2--d2，GPT对应a3--d3。四列分别为N=4、6、8和12。灰色圆点与蓝色方块分别表示关闭和开启推理。点为全部11个R值下的全过程条件均值，R按真实数值间距排列。阴影为原有的10,000次完整轨迹bootstrap逐点95\%百分位区间。竖虚线标记R=1。图中包含264个条件和6,215个完整群体，每条件10—30组。全部博弈采用无沟通、T=10、E=20、固定成员、全部历史和汇总反馈。配置配对见补充材料S1.1节。}
\label{fig:7}
\end{figure}

\begin{figure}[H]
\centering
\includegraphics[width=\linewidth]{figures/Figure_S1.pdf}
\caption{\textbf{不同结果窗口下的沟通差值。} 每个面板展示一个模型的三类差值，分别覆盖全过程、末期和终轮贡献率。估计值在N=6、T=10、E=20下对R=0.60、0.90和0.98等权平均。蓝色圆点表示首轮沟通减去无沟通，橙色方块表示每轮沟通减去无沟通，灰色三角形表示每轮沟通减去首轮沟通。误差棒保留在条件内对完整群体进行10,000次bootstrap重抽样得到的逐点95\%百分位区间。各沟通组独立抽样，相关差值复用相同的对照组抽样结果。Whole为第1--10轮，Late为第8--10轮，Final为第10轮。虚线连接彼此重叠的汇总窗口，不表示逐轮轨迹。每个窗口内的三个估计值略作横向错开。差值单位为百分点。}
\label{fig:S1}
\end{figure}

\begin{figure}[H]
\centering
\includegraphics[width=\linewidth]{figures/Figure_S2.pdf}
\caption{\textbf{沟通条件下完整群体结果的分布。} 每个模型区域包含左侧的频数带和右侧的经验累积分布，展示R=0.90、N=6、T=10、E=20时的全过程贡献率。灰色、蓝色和橙色分别表示无沟通（None）、首轮沟通（First）和每轮沟通（Every）。频数标记按五个百分点分箱，面积与群体数成比例。群体数为三或以上时标注数量。分箱左闭右开，最后一箱包含100\%。相邻标记沿纵向错开。各频数带下方的黑色标记给出由精确结果计算的中位数和四分位距。阶梯曲线使用未分箱的精确结果，表示贡献率不高于横轴数值的群体百分比。图中包含全部483个群体。DeepSeek在None、First和Every下分别包含13、10和10个群体，其余模型每种安排均包含30个群体。分布未经平滑处理。}
\label{fig:S2}
\end{figure}

\begin{figure}[H]
\centering
\includegraphics[width=\linewidth]{figures/Figure_S3.pdf}
\caption{\textbf{DeepSeek对任务设置变化的响应。} 五列分别对应R=0.60、0.90、0.98、1.00和1.02。面板a--e展示改变预告轮数（阴影行）及反馈与历史（其余行）时的全过程贡献率。轮数为T=1、5、10、20和50。No feedback/history表示同时移除反馈与历史。Aggregate和Individual分别表示汇总反馈与个体反馈，1、3或all表示保留的历史轮数。面板f--j展示改变禀赋、贡献精度和成员安排时的匹配全过程差值，对应基准分别为E=20、step/E=0.05和固定成员。点和误差棒保留全部原始均值，以及通过10,000次完整轨迹bootstrap重抽样获得的逐点95\%百分位区间。绝对贡献率与差值分别使用各自共享的横轴尺度。全部条件均采用N=6且无沟通，除被改变的参数外，E=20、T=10。随机重组使用24名参与者组成的池，其完整轨迹为独立重抽样单位。}
\label{fig:S3}
\end{figure}
